%! Unit 1: Asking Questions with Data: Study Design — Quiz 1 (Lessons 1.0–1.3) — blank
% THE REAL QUIZ. Parallel form of unit01/quiz01/sample_quiz: same blueprint and same ideas,
% different numbers and contexts, reshuffled vocabulary letters, different correct MC letters.
% Structure is deliberately identical to the sample quiz file — same parts, same guards, same
% slot order — so its pagination behaves the same way.
%
% Standalone assessment on the unit-test model (LESSON_SHAPE.md section 6): 10pt,
% \parthead{Part …} strips, \workrowsep 5pt identical in blank and key, \namedateperiod kept
% because it is taken in a testing setting. Merged into no packet. THIS IS THE GRADED FORM —
% it is built by unit01/quiz01/Makefile like the others, but it must never be handed out as
% study material the way sample_quiz is.
%
% BLUEPRINT — 19 items, 50 points, Lessons 1.0–1.3 only:
%   A vocabulary matching  8 items ×1 =  8   (1.0 ×2, 1.1 ×1, 1.2 ×3, 1.3 ×2)
%   B multiple choice      6 items ×2 = 12   (1.0, 1.1, 1.2 ×2, 1.3 ×2)
%   C short answer         4 items ×5 = 20   (C1 1.0/1.1 · C2 1.1 · C3 1.2 · C4 1.3)
%   D extended response    1 item ×10 = 10   (1.3 crux + 1.2 parameter/statistic)
% Contexts are new again: Westbrook High School (840) and Clearwater Aquatic Center (1,400).
% Riverbend, Lakeside, Harbor Point, Northgate, Cedar Ridge, and Bayside are spent in the
% lessons and the unit test; Millbrook and Stonebridge are spent in the sample quiz.
% Every number verified in pure Python before authoring.
%
% THE KEY IS GENERATED, NEVER HAND-EDITED:
%   python3 templates/lesson/mkkey.py unit01/quiz01/actual_quiz/main.tex \
%       unit01/quiz01/actual_quiz_key/main.tex unit01/quiz01/actual_quiz_key/answers.json
\documentclass[10pt]{article}
\usepackage{saar-article}
\usepackage{saar-boxes}

% Work blocks on a quiz need handwriting room, not typeset spacing. This moves the blank and
% the key together, so it cannot open a page mismatch.
\setlength{\workrowsep}{5pt}

% Part-divider strip
\newcommand{\parthead}[1]{%
  \vspace{6pt}\noindent
  \begin{headlinebox}{cerulean}{\color{white}\bfseries #1}\end{headlinebox}%
  \vspace{4pt}}

\begin{document}

\pageheader{Unit 1: Asking Questions with Data: Study Design}{Quiz 1 --- Lessons 1.0--1.3}
\namedateperiod

\vspace{0.10in}
\begin{remindbox}
\small
\textbf{Quiz 1 covers Lessons 1.0--1.3} --- data and variables, the data cycle, populations and
samples, and the four sampling techniques. \textbf{No notes and no calculator sharing.} Show
your work where a computation is asked for; an answer with no work earns no credit. Every
answer that interprets a number must name the group it describes.
\textbf{50 points:} Part A $8$, Part B $12$, Part C $20$, Part D $10$.
\end{remindbox}

% ======================================================================
\parthead{Part A --- Vocabulary \hfill 8 questions, 1 point each}

\small
Write the letter of the matching description next to each term. Each letter is used once.

\vspace{2pt}
\begin{minipage}[t]{0.36\linewidth}
\begin{enumerate}[leftmargin=1.9em, itemsep=3.5pt, topsep=1pt]
  \item \blank{0.8cm} Individual
  \item \blank{0.8cm} Quantitative variable
  \item \blank{0.8cm} Statistical question
  \item \blank{0.8cm} Sample
  \item \blank{0.8cm} Parameter
  \item \blank{0.8cm} Statistic
  \item \blank{0.8cm} Sampling frame
  \item \blank{0.8cm} Stratified sample
\end{enumerate}
\end{minipage}\hfill
\begin{minipage}[t]{0.61\linewidth}
\begin{enumerate}[label=\Alph*., leftmargin=1.9em, itemsep=3.5pt, topsep=1pt]
  \item A number computed from the sample. It changes from sample to sample.
  \item The part of the population data were actually collected from.
  \item A variable that records a measured amount with units, summarized with a mean.
  \item The list of individuals a sample can actually be drawn from.
  \item A question whose answers are expected to vary from individual to individual.
  \item The population is split into groups that are alike inside, and some are chosen at
        random from every group.
  \item One person or object that data are collected about --- one row of a data table.
  \item A number that describes the population. Usually unknown.
\end{enumerate}
\end{minipage}
\normalsize

% ======================================================================
\vspace{4pt}
\boxguard[8]
\parthead{Part B --- Multiple Choice \hfill 6 questions, 2 points each}

\small
\begin{enumerate}[leftmargin=2.2em, itemsep=5pt, topsep=2pt, start=9]

\item Westbrook High School records each student's \textbf{locker number}. This variable is
\begin{enumerate}[label=(\Alph*), leftmargin=1.8em, itemsep=1pt, topsep=1pt]
  \item categorical, because the numbers are labels for lockers.
  \item quantitative, because the values are numbers.
  \item quantitative, because the lockers can be put in order.
  \item neither, because a locker number is not really data.
\end{enumerate}
\hfill \textbf{Answer:} \blank{0.9cm}

\item Which of these is a \textbf{statistical question}?
\begin{enumerate}[label=(\Alph*), leftmargin=1.8em, itemsep=1pt, topsep=1pt]
  \item How many seats are in the Westbrook auditorium?
  \item Did the Westbrook team win its last game?
  \item How many hours do Westbrook students sleep on a school night?
  \item How many hours did I sleep last night?
\end{enumerate}
\hfill \textbf{Answer:} \blank{0.9cm}

\item Clearwater Aquatic Center's own membership records show that \textbf{exactly $30\%$ of
all $1{,}400$ members} take a class. That $30\%$ is a
\begin{enumerate}[label=(\Alph*), leftmargin=1.8em, itemsep=1pt, topsep=1pt]
  \item statistic, because it is written as a percent.
  \item statistic, because the center had to look it up.
  \item parameter, because the number came out round.
  \item parameter, because it describes every member of the population.
\end{enumerate}
\hfill \textbf{Answer:} \blank{0.9cm}

\item Two different random samples of $70$ Clearwater members give $40\%$ and $34\%$ taking a
class. This result shows
\begin{enumerate}[label=(\Alph*), leftmargin=1.8em, itemsep=1pt, topsep=1pt]
  \item sampling variability --- the parameter did not move, and both are honest estimates.
  \item that one of the two samples must have been collected incorrectly.
  \item that the parameter changed between the two surveys.
  \item that the true value must be exactly $37\%$.
\end{enumerate}
\hfill \textbf{Answer:} \blank{0.9cm}

\item A researcher stands at the Clearwater front desk on a Saturday morning and surveys
whoever walks in. His sample is
\begin{enumerate}[label=(\Alph*), leftmargin=1.8em, itemsep=1pt, topsep=1pt]
  \item a cluster sample, because he surveyed one whole morning.
  \item a convenience sample, because a member who never comes on Saturday could not be chosen.
  \item a systematic sample, because members arrive in order.
  \item a simple random sample, because he did not pick favorites.
\end{enumerate}
\hfill \textbf{Answer:} \blank{0.9cm}

\item Which sentence states the difference between a \textbf{stratum} and a \textbf{cluster}?
\begin{enumerate}[label=(\Alph*), leftmargin=1.8em, itemsep=1pt, topsep=1pt]
  \item A stratum is chosen by chance; a cluster is chosen by convenience.
  \item There is no difference --- both mean a group inside the population.
  \item A stratum is alike inside and differs from other strata; a cluster is meant to be a
        small mix of the whole.
  \item A stratum is a small mix of the whole; a cluster is alike inside.
\end{enumerate}
\hfill \textbf{Answer:} \blank{0.9cm}

\end{enumerate}
\normalsize

% ======================================================================
\newpage
\parthead{Part C --- Short Answer \& Computation \hfill 4 questions, 5 points each}

\small
\begin{enumerate}[leftmargin=2.2em, itemsep=9pt, topsep=2pt, start=15]

% ---------------------------------------------------------------- C1 (1.0 / 1.1)
\item \textbf{Westbrook High School} has $840$ students. For each student the office records
locker number, minutes spent at practice yesterday, whether the student takes the bus, and
preferred field-trip destination.

\vspace{2pt}
(a) Who are the \textbf{individuals} in this data set? \blank{6.2cm}

\vspace{3pt}
(b) Classify each variable.

\vspace{2pt}
\renewcommand{\arraystretch}{1.25}
\begin{tabularx}{\linewidth}{@{} X c @{}}
\rowcolor{frost}
\textbf{Variable} & \textbf{Categorical or quantitative?} \\
\midrule
Locker number & \blank{4.4cm} \\
Minutes at practice yesterday & \blank{4.4cm} \\
Takes the bus (yes / no) & \blank{4.4cm} \\
Preferred field-trip destination & \blank{4.4cm} \\
\end{tabularx}

\vspace{4pt}
(c) The office's first draft question is \emph{``How long is practice?''} Explain why that is
not a statistical question, then rewrite it so that it is.

\par\writeline
\par\writeline

% ---------------------------------------------------------------- C2 (1.1)
\boxguard[20]
\item \textbf{Westbrook's student council} asks $50$ students to name a preferred field-trip
destination. The counts are below.

\vspace{2pt}
(a) Complete the relative frequency column as a percent.

\vspace{2pt}
\renewcommand{\arraystretch}{1.25}
\begin{tabularx}{\linewidth}{@{} X c c @{}}
\rowcolor{frost}
\textbf{Destination} & \textbf{Frequency} & \textbf{Relative frequency} \\
\midrule
Science museum & $22$ & \blank{2.6cm} \\
Aquarium       & $14$ & \blank{2.6cm} \\
State capitol  & $9$  & \blank{2.6cm} \\
Theater        & $5$  & \blank{2.6cm} \\
\end{tabularx}

\vspace{4pt}
(b) Write a \textbf{conclusion in context} about the most popular destination, naming the group
it applies to.

\par\writeline
\par\writeline

\vspace{2pt}
(c) The council writes: \emph{``Most Westbrook students want the science museum.''} Give the
two things that are wrong with that sentence.

\par\writeline
\par\writeline

% ---------------------------------------------------------------- C3 (1.2)
\boxguard[18]
\item \textbf{Clearwater Aquatic Center} has $1{,}400$ members. The manager surveys $70$
members chosen at random; $28$ of them take a class.

\vspace{2pt}
(a) Population: \blank{4.4cm} \quad Sample size: \blank{1.8cm}

\vspace{3pt}
Sample: \blank{9.0cm}

\vspace{3pt}
(b) What percent of the sample takes a class, and about how many of all $1{,}400$ members does
that estimate?

\begin{work}
  \dfrac{28}{70} &= 0.40 = 40\% \\
  0.40 \times 1400 &= 560 \text{ members}
\end{work}

(c) Is the $40\%$ a parameter or a statistic? Say how you know.

\par\writeline

(d) Name one \textbf{constraint} that keeps the manager from taking a census of all $1{,}400$
members.

\par\writeline

% ---------------------------------------------------------------- C4 (1.3)
\newpage
\item \textbf{Westbrook High School} has $840$ students, a numbered list of all of them, and
$30$ homerooms of $28$ students each.

\vspace{2pt}
(a) Name the sampling technique used in each row.

\vspace{2pt}
\renewcommand{\arraystretch}{1.3}
\begin{tabularx}{\linewidth}{@{} X c @{}}
\rowcolor{frost}
\textbf{What the researcher did} & \textbf{Technique} \\
\midrule
Put all $840$ student numbers in a generator and drew $42$. & \blank{3.2cm} \\
Split the school by grade and drew students at random from every grade, in proportion. &
  \blank{3.2cm} \\
Ordered the list and took every $20$th student after a random start. & \blank{3.2cm} \\
Drew $2$ homerooms at random and surveyed everyone in them. & \blank{3.2cm} \\
\end{tabularx}

\vspace{4pt}
(b) For a \textbf{systematic} sample of $42$ students, show why $k = 20$, then list the first
four students chosen from a random start of $6$.

\begin{work}
  k &= N \div n = 840 \div 42 = 20
\end{work}

\vspace{-2pt}
\blank{1.6cm} \quad \blank{1.6cm} \quad \blank{1.6cm} \quad \blank{1.6cm}

\vspace{4pt}
(c) For a \textbf{simple random} sample, a generator gives $517$, $926$, $304$, $688$, $150$.
One of those numbers cannot be used. Which one, and why?

\par\writeline

\end{enumerate}
\normalsize

% ======================================================================
\vspace{4pt}
\boxguard[10]
\parthead{Part D --- Extended Response \hfill 1 question, 10 points}

\small
\begin{enumerate}[leftmargin=2.2em, itemsep=10pt, topsep=2pt, start=19]

\item \textbf{Westbrook} ($840$ students) wants to know whether students would use a Saturday
tutoring session. The office has a numbered list of all $840$ students, the four grades are
believed to differ a lot in their answers, and the $30$ homerooms are \textbf{built by grade}.

\vspace{2pt}
(a) Choose \textbf{one} of the four sampling techniques and justify it from what this context
allows. Your reason must name \emph{how the groups are built}.

\par\writeline
\par\writeline

\vspace{2pt}
(b) Instead, a teacher draws $2$ homerooms at random and surveys all $56$ students in them;
$42$ say yes. Show her two computations: the percent of her sample saying yes, and the number
of all $840$ students that percent points to.

\begin{work}
  \dfrac{42}{56} &= 0.75 = 75\% \\
  0.75 \times 840 &= 630 \text{ students}
\end{work}

(c) Her arithmetic is correct. Explain what is wrong with her report anyway, and state what she
\emph{is} allowed to say.

\par\writeline
\par\writeline
\par\writeline

\vspace{2pt}
(d) Is her $75\%$ a parameter or a statistic? Name one change to her method that would make it
an estimate worth trusting.

\par\writeline
\par\writeline

\end{enumerate}
\normalsize

\end{document}
